\documentclass[11pt]{article}
\pdfminorversion=7
\pdfobjcompresslevel=0

% Change "review" to "final" to generate the final version.
% Change to "preprint" to generate a non-anonymous version with page numbers.
\usepackage[preprint]{acl}

% Standard package includes
\usepackage{times}
\usepackage{latexsym}
\usepackage[T1]{fontenc}
\usepackage[utf8]{inputenc}
\usepackage{microtype}
\usepackage{inconsolata}
\usepackage{graphicx}
\usepackage{booktabs}
\usepackage{amsmath}
\usepackage{amssymb}
\usepackage{xcolor}

\title{Crayotter: Artifact-Grounded Multimodal Agents for Long-Form Video Editing}

\author{
  Lecheng Yan$^{*\dagger}$,
  Yichong Zhang$^{*\dagger}$,
  Ben Pan$^{*}$,
  Xiaoyu Zheng$^{*}$,
  Jiawei Qian$^{*}$,
  Anqi Wu$^{*}$, \\
  {\bf Wenxi Li},
  {\bf Chenyang Lyu} \\
  Crayotter Project
}

\begin{document}
\maketitle

\begingroup
\renewcommand{\thefootnote}{\fnsymbol{footnote}}
\footnotetext[1]{Core Contributor.}
\footnotetext[2]{Project Lead.}
\endgroup

\begin{figure*}[t]
  \centering
  \includegraphics[width=0.98\textwidth]{fig/crayottor_visual.pdf}
  \caption{Crayotter client and case-study editing trajectory. Left: the
  workbench entry exposes task history, local assets, and the fallback path to
  agent-searched external assets. Center: a ReAct editing timeline for a
  travel-video request, where the agent performs precision cuts, arranges
  cross-dissolve transitions, and aligns narration and subtitle tracks with
  clip-level content. Right: multimodal analysis, reflection, and optional
  human review form an explicit revise loop before final export.}
  \label{fig:analysis}
\end{figure*}


\begin{abstract}
Recent progress in LLM agents has improved short-horizon tool use, but long-horizon multimodal editing remains difficult: decisions are weakly grounded, intermediate states are hard to audit, and failures are expensive to localize. We formulate long-form video editing as an \emph{artifact-grounded agent trajectory} problem, where planning, execution, and revision are all conditioned on explicit external states. We present \textbf{Crayotter}, an open-source multimodal multi-agent framework that maps one natural-language intent to a complete edited video via three phases: material preparation, deep editing research, and tool-grounded agent execution. Our central mechanism is \textbf{environment-grounded reflection}: the agent updates its policy from observable artifacts instead of relying only on latent reasoning traces. We further propose a trajectory-level RLVR framework with a three-tier cascade reward that combines verifiable editing signals, LLM-as-judge multi-dimensional scoring, and human preference calibration, optimized via GRPO for sample-efficient policy improvement. Code and examples are publicly available at
\url{https://github.com/idwts/Crayotter}.
\end{abstract}

\section{Introduction}

Long-form visual storytelling is rapidly shifting from single-model generation
to production-oriented agent systems. Large language models now support
multi-step reasoning, tool use, and role-based collaboration
\cite{react23,autogen23,metagpt23}. In video creation, recent systems decompose
production into explicit stages such as story planning, shot generation,
review, cinematography, and post-production
\cite{videoDirectorGPT24,dreamFactory24,movieAgent25,mmStoryAgent25,filmAgent25,
animaker25}. FilmAgent, for example, makes virtual filmmaking concrete by
simulating crew roles in 3D spaces, while AniMaker shows how multi-agent
animation can benefit from candidate exploration and storytelling-aware
evaluation. A common lesson is emerging: high-quality AI video depends not only
on stronger generators, but also on controllable production workflows.

Long-form video \emph{editing}, however, remains under-formalized from an
agent-systems perspective. Editing is not simply a matter of prompting a video
generator. It requires a sequence of concrete decisions over heterogeneous
artifacts: user intent, local or retrieved footage, multimodal clip analyses,
timeline structures, transitions, narration, subtitles, audio levels, tool logs,
and rendered feedback. A small error in asset selection or clip timing can
propagate into broken continuity, unnatural pacing, or narration that no longer
matches the image. This makes editing a long-horizon, multimodal control problem
where state grounding, trajectory reproducibility, and failure diagnosis are
central.

Figure~\ref{fig:analysis} illustrates the kind of workflow we target. A user
starts from a high-level travel-video request, the system prepares or searches
for visual material, and the agent exposes a ReAct-style editing trajectory
rather than a black-box answer. Cuts, transitions, narration, subtitles, and
rendered previews become inspectable units. When the result fails, the agent can
revise the affected timeline segment or hand control to a human editor instead
of restarting the entire production. This is the core intuition behind
Crayotter: a video-editing agent should make its production state visible,
reusable, and correctable.

Existing work only partially covers this setting. Generation-first pipelines are
strong at narrative synthesis, and recent long-context models improve temporal
consistency beyond single clips \cite{fifo24,lct25,mint25,oneMinute25}. Yet
these methods usually focus on synthesizing shots rather than coordinating
editing-time operations over real footage. Tooling-oriented projects improve
practical usability, but they often lack scientific trajectory logs that allow a
run to be replayed, diagnosed, or compared. Meanwhile, evaluation remains
fragmented for editing-specific qualities such as cut continuity, transition
naturalness, narration alignment, and localized error attribution
\cite{dreamStory24,storyIter24,viStoryBench26,evalCrafter24}.

We present \textbf{Crayotter}, a multimodal multi-agent framework for
prompt-driven video editing over heterogeneous source material. Given a user
intent, Crayotter executes three phases. In \emph{material preparation}, it
expands the request into coverage tags, retrieves or uses local videos, reranks
candidate assets, and produces multimodal analysis artifacts. In \emph{deep
editing research}, it synthesizes those artifacts into an editing blueprint,
including narrative structure, clip order, rhythm, transitions, and narration.
In \emph{tool-grounded agent execution}, it translates the blueprint into
editing actions such as cutting, timeline construction, transition insertion,
subtitle generation, loudness adjustment, and export.

The key mechanism is \textbf{environment-grounded reflection}. Rather than
treating an LLM conversation as the only state, Crayotter externalizes state
into artifacts that can be inspected by both agents and users: retrieval
coverage reports, analysis JSON, timeline plans, transition plans, tool calls,
intermediate renders, and final outputs. Inspired by process-oriented research
agents and artifact-based reflection \cite{researStudio25,deepPresenter26}, the
system revises from observable environment feedback and selectively re-executes
failed segments. This design turns long-form editing from a single opaque
generation attempt into a reproducible trajectory of planning, execution, and
repair.

\textbf{Our contributions are four-fold.}
\begin{enumerate}
\item We formulate long-form multimodal video editing as an
artifact-grounded agent trajectory problem, where source assets, plans, timeline
states, renders, and logs are explicit state variables for reflection and
repair.
\item We introduce a coverage-aware multimodal footage retrieval loop that
decomposes an editing request into visual, narrative, style, and action
coverage tags, reranks candidate videos with sampled frames and temporal
windows, and iteratively searches for missing semantic evidence.
\item We deliver Crayotter as an open-source three-phase editing system with
structured editing research, tool-grounded execution, runtime observability,
configurable phase routing, an experience memory system, and a workbench client
for human-agent partnership.
\item We propose a trajectory-level RLVR framework using GRPO with a three-tier
cascade reward that combines verifiable tool signals, episode-level outcome
scoring, and human preference calibration.
\end{enumerate}

\begin{figure*}[t]
  \centering
  \includegraphics[width=\textwidth]{fig/crayottor_framework.pdf}
  \caption{Overall architecture of Crayotter. The system follows a three-phase
  pipeline (material preparation, editing research, and tool-grounded agent execution)
  with artifact-grounded reflection, runtime observability, and reusable
  trajectory logging.}
  \label{fig:overview}
\end{figure*}

\section{Related Work}

\subsection{Agentic Story and Film Generation}

Multi-agent systems have recently been used to automate story and film
production. VideoDirectorGPT introduces LLM-guided planning for multi-scene
video generation with layout control \cite{videoDirectorGPT24}. DreamFactory and
MovieAgent further push multi-agent coordination for long-video production
\cite{dreamFactory24,movieAgent25}. MM-StoryAgent extends the pipeline to
cross-modal text-image-audio storybooks \cite{mmStoryAgent25}, while FilmAgent
targets end-to-end virtual 3D film automation \cite{filmAgent25}. AniMaker adds
MCTS-guided candidate exploration and animation-specific evaluation
(AniEval) for multi-shot narrative animation \cite{animaker25}.

Compared to these systems, Crayotter focuses on practical editing over
retrieved real videos, emphasizing artifact-observable execution, reproducible
tool traces, and post-production integration rather than generation-first
virtual production.

\subsection{Long-Horizon Video Editing and Cinematic Compilation}

Several recent systems are closer to Crayotter because they operate on existing
or long-form video materials rather than only synthesizing new clips. DIRECT
formulates video mashup creation as hierarchical multi-agent planning, with
screenwriter, director, and editor roles coordinating structural anchoring,
editing intent, and fine-grained shot-sequence optimization
\cite{li2026direct}. CineAgents targets instruction-driven cinematic video
compilation, using CineBench for evaluation and a design-and-compose workflow
with hierarchical narrative memory and iterative narrative planning
\cite{zhang2026benchmark}. CutClaw addresses hours-long video editing under
music synchronization, combining multimodal decomposition, long-range planning,
and fine-grained clip selection \cite{Zhao2026CutClaw}.

Open-source editing tools also show strong practical momentum.
FireRed-OpenStoryline supports conversational video creation with online media
search, LLM-powered planning, precise tool orchestration, human-in-the-loop
control, and reusable style skills.\footnote{\url{https://github.com/FireRedTeam/FireRed-OpenStoryline}} NarratoAI
focuses on one-click film commentary and automatic editing workflows, providing
a product-oriented pipeline for script generation, narration, and video
assembly.\footnote{\url{https://github.com/linyqh/NarratoAI}} These projects demonstrate usable interfaces and
task-specific automation. Crayotter differs by making the editing process a
research object: it externalizes retrieval evidence, analysis artifacts,
timeline states, tool logs, and rendered feedback so that a run can be replayed,
diagnosed, selectively repaired, and eventually optimized with trajectory-level
signals.

Evaluation is increasingly central for story visualization and long-form video.
DreamStory introduces multi-subject consistency modeling and DS-500 benchmark
\cite{dreamStory24}. Story-Iter proposes iterative long-story refinement
\cite{storyIter24}. ViStoryBench provides comprehensive story visualization
metrics and human-verified scripts \cite{viStoryBench26}. EvalCrafter analyzes
evaluation reliability for video generation models \cite{evalCrafter24}.

Crayotter can reuse such benchmarks while adding editing-specific axes such as
cut continuity, transition naturalness, and narration timeline validity.

\subsection{Long-Context and Multi-Event Video Modeling}

Long context remains a bottleneck in video generation. FIFO-Diffusion explores
infinite-length inference without retraining \cite{fifo24}. Long Context Tuning
expands context windows of pre-trained single-shot diffusion models for coherent
multi-shot scenes \cite{lct25}. MinT introduces temporal control by binding
events to time windows \cite{mint25}. One-Minute Video Generation with test-time
training studies longer narrative generation via architectural adaptation
\cite{oneMinute25}.

These approaches focus on generative backbones, while our system-level setting
operates on editing-time composition and revision of heterogeneous source
materials.

\subsection{Deep Research Training and Reflection}

DR Tulu proposes reinforcement learning with evolving rubrics for long-form deep
research tasks \cite{drTulu25}; SAGE diagnoses retrieval bottlenecks in deep
research agents \cite{sage26}. DeepPresenter shows that environment-grounded
reflection over rendered artifacts can outperform fixed template workflows in
presentation generation \cite{deepPresenter26}. This directly motivates our
editing reflection loop over rendered video artifacts.

\section{Method: Crayotter Architecture}

Crayotter maps an open-ended editing request to a finished video through a
sequence of inspectable intermediate representations rather than through a
single latent generation step. Given a user request $q$ and an optional local
asset set $\mathcal{A}_0$, the system realizes the transformation
\[
    \mathcal{F}: (q,\mathcal{A}_0) \rightarrow y,
\]
where $y$ is the exported video. We factorize $\mathcal{F}$ into three
production phases:
\[
    (q,\mathcal{A}_0)
    \xrightarrow{\text{retrieval}}
    \mathcal{P}
    \xrightarrow{\text{research}}
    \mathcal{B}
    \xrightarrow{\text{tools}}
    y.
\]
Here $\mathcal{P}$ is a coverage-verified material pool, $\mathcal{B}$ is a
time-grounded editing blueprint, and the final phase executes $\mathcal{B}$ in
an external editing environment. This decomposition follows the production
logic of recent agentic visual systems: every phase produces artifacts that are
visible to later agents, to the tool environment, and to a human reviewer.

Figure~\ref{fig:ui} shows the corresponding workbench interface through which
users inspect tasks, assets, and intermediate agent outputs during execution.

\begin{figure}[t]
  \centering
  \includegraphics[width=\columnwidth]{fig/crayottor_ui.pdf}
  \caption{Crayotter workbench interface. The client exposes task history,
  local asset slots, execution status, and agent-generated intermediate
  artifacts in one workspace, supporting human-agent inspection during
  long-horizon editing.}
  \label{fig:ui}
\end{figure}

\begin{figure*}[t]
  \centering
  \includegraphics[width=\textwidth]{fig/crayottor_search.pdf}
  \caption{Coverage-aware multimodal footage retrieval. Crayotter converts an
  abstract editing request into concrete coverage tags, retrieves a high-recall
  video pool, reranks candidates through frame and temporal-window analysis,
  verifies top-$K$ videos with full-video understanding, and turns uncovered
  tags into follow-up searches until the material pool is sufficient or the
  retrieval budget is exhausted.}
  \label{fig:retrieval}
\end{figure*}

\subsection{Task Formulation and Intermediate Artifacts}

We define a source video as $v_j=(x_j,a_j,L_j)$, with visual stream $x_j$,
optional audio stream $a_j$, and duration $L_j$. An edited video is represented
as a timeline
\[
    y=\mathrm{Render}(\mathcal{E}), \quad
    \mathcal{E}=\{e_\ell\}_{\ell=1}^{N},
\]
where each event $e_\ell=(v_{j_\ell},\tau_\ell,o_\ell)$ selects a source
video, a temporal interval
$\tau_\ell=(t^{\mathrm{in}}_\ell,t^{\mathrm{out}}_\ell)$, and an editing
operator $o_\ell$ such as cut, transition, narration alignment, subtitle
placement, loudness normalization, or export. The core difficulty is that
$\tau_\ell$ and $o_\ell$ must be chosen from imperfectly retrieved material
while remaining faithful to the global request $q$.

Crayotter therefore introduces three explicit artifact types:
\[
    \begin{aligned}
    \mathcal{P} &= (\mathcal{S},\Gamma,\mathcal{M}),\\
    \mathcal{B} &= \{b_\ell\}_{\ell=1}^{N},\\
    \mathcal{H} &= \{(s_k,a_k,s_{k+1},\rho_k)\}_{k=1}^{K}.
    \end{aligned}
\]
$\mathcal{S}$ is the selected source pool, $\Gamma$ stores coverage evidence,
$\mathcal{M}$ stores multimodal analysis records, $\mathcal{B}$ is a
shot-level editing blueprint, and $\mathcal{H}$ is the tool execution history
with observable rewards or diagnostics $\rho_k$. The following subsections
describe how each artifact is constructed.

\subsection{Phase 1: Coverage-Aware Multimodal Footage Retrieval}

Source-footage preparation is a central bottleneck in practical video editing:
the editor cannot plan a convincing timeline unless the material pool contains
visual evidence for the requested story. We formulate retrieval as weighted tag
coverage. The request $q$ is expanded into production-oriented tags
\[
    T = T^{\mathrm{scene}} \cup T^{\mathrm{act}}
      \cup T^{\mathrm{style}} \cup T^{\mathrm{story}}
      \cup T^{\mathrm{shot}},
\]
where the dimensions encode scene requirements, people/actions, affective
style, narrative purpose, and expected shot evidence. Each tag $t_i$ is assigned
an importance weight $w_i$.

At retrieval round $r$, a query generator produces $q_r$ from the original
request and the currently missing tags. A coarse retriever returns a high-recall
candidate list
\[
    \begin{aligned}
    \mathcal{C}_r =
    \operatorname{TopM}_{v\in\mathcal{D}}
    \bigl[&
    \alpha\,s_{\mathrm{text}}(q_r,v)\\
    &+(1-\alpha)\,s_{\mathrm{meta}}(q_r,v)
    \bigr],
    \end{aligned}
\]
where $\mathcal{D}$ is an external video index. For each candidate $v$, the
system samples frames $F(v)$ and temporal windows $W(v)$, then estimates a
tag-support vector
\[
    g(v,t_i)=
    \max_{u\in F(v)\cup W(v)}
    p_{\psi}(t_i \mid u,q),
\]
where $p_{\psi}$ denotes the multimodal understanding module. Candidates are
reranked by request relevance and marginal coverage gain:
\[
    \begin{aligned}
    D_i(v,\mathcal{S})
    &= \max\bigl(0,g(v,t_i)-G_{\mathcal{S}}(t_i)\bigr),\\
    \operatorname{score}(v\mid\mathcal{S},T)
    &= \lambda\,s_{\mathrm{rel}}(q,v)\\
    &\quad +(1-\lambda)\sum_i w_i\,D_i(v,\mathcal{S}),
    \end{aligned}
\]
with $G_{\mathcal{S}}(t_i)=\max_{v\in\mathcal{S}}g(v,t_i)$. The top-$K$
videos are then downloaded for full-video verification. The selected pool
$\mathcal{S}$ is evaluated by the normalized coverage objective
\[
    \operatorname{Cov}(\mathcal{S},T)
    =\frac{\sum_i w_i\,\mathbb{I}[G_{\mathcal{S}}(t_i)\ge \eta_i]}
           {\sum_i w_i}.
\]
Uncovered tags form $\Delta_r=\{t_i\mid G_{\mathcal{S}}(t_i)<\eta_i\}$ and are
converted into the next query $q_{r+1}$. The loop stops when
$\operatorname{Cov}(\mathcal{S},T)$ reaches a target threshold or the retrieval
budget is exhausted. Phase~1 outputs the material pool
$\mathcal{P}=(\mathcal{S},\Gamma,\mathcal{M})$, where $\Gamma$ records
tag-level evidence and $\mathcal{M}$ stores per-video multimodal analyses.

\subsection{Phase 2: Deep Editing Research}

Phase~2 transforms the material pool into an editing blueprint without
executing editing tools. Its key problem is temporal localization: the agent
must decide not only which clip to use, but which in/out point supports a
narrative beat. We borrow the interface principle from our timestamp-watermark
technical report: many frozen MLLMs localize events more reliably when temporal
coordinates are rendered directly on the perceptual evidence. For each selected
video $v_j$, Crayotter constructs a time-marked view
\[
    \tilde{v}_j = R_{\theta}(v_j),
\]
where $R_{\theta}$ overlays human-readable time coordinates on sampled frames
or analysis windows. The model parameters are unchanged; only the input
interface is modified. Because the timestamp is drawn on the same frame that
contains the event evidence, the research agent can bind semantic observations
to absolute editing coordinates.

Let $\mathcal{Z}=\{z_m\}$ be narrative beats derived from $q$ and the retrieval
report, such as establishing shot, activity montage, close-up interaction, or
closing scene. For each beat, the research agent predicts a source and interval
\[
    (\hat{j}_m,\hat{\tau}_m)
    = \operatorname{parse}
      \bigl(f_{\phi}(\{\tilde{v}_j\}_{v_j\in\mathcal{S}},z_m,\Gamma)\bigr),
\]
then assembles a blueprint
\[
    b_m=(v_{\hat{j}_m},\hat{\tau}_m,c_m,r_m,n_m),
\]
where $c_m$ is the intended cinematic role, $r_m$ is the transition or rhythm
rationale, and $n_m$ is narration or subtitle intent. The result
$\mathcal{B}=\{b_m\}$ is therefore not free-form prose: it is a temporally
addressable production plan whose intervals can be inspected, revised, and
executed by tools. This design localizes later failures to a specific source
clip, timestamp span, or planning rationale.

\subsection{Phase 3: Tool-Grounded Timeline Execution}

Phase~3 realizes the blueprint through external editing tools. We model the
editing environment as an artifact transition system. At step $k$, the state
$s_k$ contains the current timeline, source assets, subtitle tracks, narration
files, rendered previews, and tool logs. The editor chooses an action
$a_k\in\mathcal{U}$ from the tool set:
\[
    s_{k+1}=E(s_k,a_k;\mathcal{B}),
\]
where $E$ is the editing environment. Typical actions include trimming,
placing a clip on the timeline, inserting a transition, aligning narration,
generating subtitles, normalizing loudness, and exporting a preview.

The key mechanism is environment-grounded reflection. After each tool call, a
verifier computes artifact-observable diagnostics
\[
    \rho_k =
    R_{\mathrm{tool}}(s_k,a_k,s_{k+1},\mathcal{B},q),
\]
including tool validity, timestamp accuracy, coverage preservation, narration
alignment, transition smoothness, and render quality. If a diagnostic fails,
the agent observes the concrete artifact that failed and repairs only the
affected segment rather than restarting the full edit.

\subsection{Agent Roles and Artifact Contracts}

The implementation follows an explicit planner-researcher-executor separation.
The planner converts user intent into phase goals, the researcher synthesizes
multimodal evidence into an editable blueprint, and the executor translates the
blueprint into concrete editing tool calls. Across all phases, state is
externalized as artifacts rather than hidden in latent traces.
A mandatory \emph{composite re-analysis} step occurs before narration: the
merged video is re-analyzed via the multimodal vision model to ensure narration
scripts align with actual rendered content rather than only source-material
analysis.

We maintain a stable artifact contract:
\begin{itemize}
\item input-side artifacts: request, retrieval candidates, and per-video
analysis outputs;
\item planning artifacts: narrative blueprint, clip ordering, rhythm targets,
transition and narration plans;
\item execution artifacts: timeline files, intermediate renders, and tool logs.
\end{itemize}
This contract allows deterministic replay, phase-level diagnosis, and selective
re-execution after failures.

\subsection{Modular Tool Taxonomy}

Crayotter exposes 21 registered tools partitioned into two disjoint sets:
Phase~1 preparation tools and Phase~3 editing tools.
Table~\ref{tab:tools} summarizes the full taxonomy.

\begin{table}[t]
\caption{Tool taxonomy with functional grouping.}
\label{tab:tools}
\centering
\small
\resizebox{\columnwidth}{!}{
\begin{tabular}{ll}
\toprule
Category & Tools \\
\midrule
Search & \texttt{search\_video} \\
Ranking & \texttt{rank\_video\_candidates} \\
Download & \texttt{download\_video} \\
Analysis & \texttt{analyze\_video}, \texttt{recall\_semantic\_segments} \\
Cutting & \texttt{cut\_video}, \texttt{batch\_cut\_video} \\
Merging & \texttt{merge\_videos} \\
Inspection & \texttt{inspect\_video\_duration} \\
Transitions & \texttt{add\_transition}, \texttt{plan\_transition\_timeline}, \\
 & \texttt{list\_transition\_presets} \\
Continuity & \texttt{score\_cut\_continuity}, \\
 & \texttt{recommend\_transition\_for\_cut} \\
Timeline & \texttt{build\_edit\_timeline\_from\_segments}, \\
 & \texttt{align\_narration\_to\_timeline}, \\
 & \texttt{validate\_timeline\_constraints} \\
Narration & \texttt{add\_narration\_segments}, \\
 & \texttt{validate\_narration\_timeline} \\
Audio Post & \texttt{duck\_background\_audio}, \texttt{normalize\_loudness} \\
Subtitles & \texttt{add\_subtitles} \\
Export & \texttt{export\_video} \\
\bottomrule
\end{tabular}
}
\end{table}

The transition subsystem supports 17 preset effects organized into three
categories---basic (crossfade, fade-through-black/white), motion (wipe and
slide in four directions), and cinematic (zoom-in, smooth-left/right,
distance)---each with configurable duration and per-cut-point assignment via
a \texttt{transition\_plan} schema. Continuity tools sample frame signatures
around each cut, estimate visual gaps from brightness, saturation, and motion,
and recommend hard cuts, dissolves, or fade-through-black transitions. These
modules are deliberately lightweight: they provide observable diagnostics for
the reflection loop without turning the method into a collection of unrelated
engineering features.

Several runtime utilities support the same artifact contract. Phase routing can
run the full pipeline, skip the research phase, operate directly on local
materials, or perform a local-first sufficiency check before online retrieval.
A bounded reference-only memory stores reusable tool patterns, failure guards,
and quality checklists while stripping task-specific creative content. Together
with an agent stall watchdog, these utilities improve replayability and
robustness but do not change the core three-phase formulation.

\subsection{Trajectory-Level RLVR}

The artifact transition view naturally defines a reinforcement learning with
verifiable rewards (RLVR) environment for Phase~3. We optimize a policy
$\pi_{\theta}$ over editing-tool actions with GRPO:
\[
    \begin{aligned}
    \max_{\theta}\;
    &\mathbb{E}_{a_k\sim\pi_{\theta}(\cdot\mid s_k,\mathcal{B})}\\
    &\left[\sum_{k=1}^{K}
    R_{\mathrm{tool}}(s_k,a_k,s_{k+1},\mathcal{B},q)\right].
    \end{aligned}
\]
Unlike preference-only rewards, $R_{\mathrm{tool}}$ is grounded in observable
editing artifacts: parsable timelines, valid file outputs, measured durations,
coverage records, preview renders, and execution logs. The reward is organized
as a three-tier cascade: step-level verifiable signals for tool execution and
artifact production, episode-level outcome scoring over export success,
duration adherence and efficiency, and human preference calibration.
The concrete shaping constants are implementation details and are listed in
Appendix~\ref{app:rlvr_rewards}.

Training episodes are specified by fixtures containing the user request, target
duration, allowed tools, seeded workspace files, and scripted warm-up turns when
needed. The implementation exports these fixtures to a \texttt{verl}-compatible
JSONL format and uses an sglang multi-turn rollout loop with subprocess-backed
tool calls. This setup lets us validate reward plumbing and trajectory replay
locally before scaling the policy-improvement run.

\section{Competitive Evaluation}

We evaluate Crayotter on a fixed pack of 23 editing themes against two
open-source or practical baselines: CapCut-Mate and CutClaw. Each generated
video is rated on a five-point scale along five dimensions: theme alignment,
content richness, narrative coherence, editing smoothness, and visual quality.
The overall score is the mean of the five dimension scores.

\begin{table*}[t]
\caption{Human evaluation by dimension. Each cell is averaged over three
annotators and 23 themes.}
\label{tab:eval_dims}
\centering
\resizebox{0.8\textwidth}{!}{
\begin{tabular}{lcccccc}
\toprule
Method & Theme & Richness & Narrative & Smoothness & Visual & Overall \\
\midrule
Crayotter (ours) & \textbf{3.59} & \textbf{3.35} & \textbf{3.22} & \textbf{3.29} & \textbf{3.54} & \textbf{3.40} \\
CapCut-Mate & 2.59 & 2.71 & 2.01 & 2.13 & 2.74 & 2.44 \\
CutClaw & 1.59 & 1.64 & 1.72 & 1.86 & 1.67 & 1.70 \\
\bottomrule
\end{tabular}
}
\end{table*}

\subsection{Scoring Protocol}

We report both AI and human judgments. For AI scoring, we use the same
multidimensional rubric across all 23 themes and all methods. For human
scoring, three annotators score the same outputs; we first average annotator
scores for each theme-method pair, then average across the 23 themes. This
case-level aggregation avoids giving extra weight to any single annotator file.

\begin{table}[t]
\caption{Overall comparison on 23 themes. Scores are mean $\pm$ standard
deviation across themes on a 1--5 scale.}
\label{tab:eval_main}
\centering
\resizebox{0.9\linewidth}{!}{
\begin{tabular}{lcc}
\toprule
Method & Human score & AI score \\
\midrule
Crayotter (ours) & $\mathbf{3.40}\pm0.59$ & $\mathbf{2.39}\pm0.62$ \\
CapCut-Mate & $2.44\pm0.72$ & $2.10\pm0.64$ \\
CutClaw & $1.70\pm0.62$ & $1.57\pm0.40$ \\
\bottomrule
\end{tabular}
}
\end{table}


The results show that Crayotter substantially outperforms both baselines under
human and AI evaluation. Human raters give Crayotter the strongest scores across
all five dimensions, which is consistent with the system's explicit material
preparation and artifact-grounded planning stages.

\section{Discussion}

\subsection{Process-Level Diagnostics}

Our analysis focuses on whether the system remains controllable during
long-horizon execution: where plans drift, which artifacts expose failure early,
and which revisions yield the highest recovery gain. This process-level view is
critical for transforming case anecdotes into actionable engineering signals.

\subsection{System Visualization and Case Showcase}

Figure~\ref{fig:analysis} is a concrete case-study view. It combines the
client-side entry panel, runtime trace surfaces, and an editing-timeline
walkthrough for a travel-video request. The left side
shows the user-facing workbench: task history, local asset slots, and the
fallback path to agent-searched external assets when local material is
insufficient. This grounds the system in an observable execution environment
instead of a prompt-only interaction.


The center panel makes the editable units explicit: raw clips, agent-guided
precision cuts, transition operators, narration planning, subtitle alignment,
and loudness adjustment all appear as first-class timeline artifacts. This is
important for our claim of environment-grounded reflection, because the agent
does not revise from hidden chain-of-thought alone; it revises from rendered
timeline states, multimodal analysis outputs, and tool-produced intermediate
artifacts that can be inspected by both the system and the user.

The right side of Figure~\ref{fig:analysis} highlights the control loop we care
about most in practice: observe, diagnose, revise, and optionally hand control
to a human editor. In our setting, debugging a failed run rarely means starting
over from scratch. Instead, users can trace failures back to specific asset
choices, narration mismatches, or timeline decisions, then revise only the
affected stage. This figure therefore serves as both a UI snapshot and a
process diagram for controllable long-horizon video editing.

\subsection{Limitations and Scope}

The 23-theme evaluation provides a competitive picture of end-to-end editing
quality across representative long-form requests. Its scope is deliberately
system-level: it measures theme alignment, content richness, narrative
coherence, editing smoothness, and visual quality, but it does not isolate the
individual contribution of every retrieval, reflection, or tool-control module.
The RLVR formulation is implemented with verifiable reward traces and a
replayable trajectory schema; large-scale policy optimization and broader human
preference studies remain outside the scope of this report.

\section{Conclusion}

This paper presented Crayotter, an open-source multimodal multi-agent framework
for long-form video editing. We formulated editing as an artifact-grounded agent
trajectory problem and showed how material preparation, deep editing research,
tool-grounded execution, and environment-grounded reflection provide explicit
states for planning, auditing, and revision. The report also introduced a
trajectory-level RLVR formulation and a 23-theme evaluation with both AI and
human scoring against CapCut-Mate and CutClaw baselines. Together, these
components show that long-form editing agents can be evaluated and improved
through observable artifacts rather than latent reasoning traces alone.

\bibliography{sample-bibliography}

\appendix

\section{RLVR Reward Details}
\label{app:rlvr_rewards}

The trajectory reward used in our implemented RLVR module combines step-level
verifiable signals and episode-level outcome scoring. At each tool call, the
agent receives a tool-success reward of $+0.6$ or a failure penalty of $-0.8$,
an artifact-production bonus of $+0.1$ when valid files are created, a non-zero
return-code penalty of $-0.1$, a repeat-call penalty of $-0.2$, and an ordering
bonus or penalty in the range $\pm0.1$--$0.4$ for procedure-sensitive actions.

At episode end, the reward adds an export bonus of $+1.0$, a duration-accuracy
term clipped between $-0.5$ and $+0.8$ using
$0.8-|d_{\text{final}}-d_{\text{target}}|/d_{\text{target}}$, a completion
bonus of $+0.3$ for non-empty final output, and an efficiency penalty of
$-0.02$ per tool call beyond six calls. The third tier supports human preference
calibration; in the experiments reported here, Tiers~1--2 provide the active
training signal while human judgments are reported as evaluation results.

\end{document}
